\begin{figure}[htbp]
    \centering
    \begin{tikzpicture}[
        node distance=0.55cm,
        block/.style={
            draw,
            rounded corners,
            align=center,
            minimum width=2.3cm,
            minimum height=0.85cm,
            font=\small
        },
        arrow/.style={-{Latex[length=2mm]}, thick}
    ]
        \node[block] (input) {Input\\$224 \times 224 \times 3$};
        \node[block, right=of input] (conv1) {Conv block 1\\$3 \rightarrow 64 \rightarrow 64$};
        \node[block, right=of conv1] (pool1) {Max pool\\$112 \times 112$};
        \node[block, right=of pool1] (conv2) {Conv block 2\\$64 \rightarrow 128 \rightarrow 128$};
        \node[block, below=of conv2] (pool2) {Max pool\\$56 \times 56$};
        \node[block, left=of pool2] (conv3) {Conv block 3\\$128 \rightarrow 256$};
        \node[block, left=of conv3] (gap) {Adaptive\\average pool};
        \node[block, left=of gap] (head) {Dropout $0.2$\\Linear $256 \rightarrow 9$};

        \draw[arrow] (input) -- (conv1);
        \draw[arrow] (conv1) -- (pool1);
        \draw[arrow] (pool1) -- (conv2);
        \draw[arrow] (conv2) -- (pool2);
        \draw[arrow] (pool2) -- (conv3);
        \draw[arrow] (conv3) -- (gap);
        \draw[arrow] (gap) -- (head);
    \end{tikzpicture}
    \caption{Architecture of the standalone $224 \times 224$ baseline CNN.
    Each convolution is followed by batch normalization and a SiLU activation.}
    \label{fig:baseline224-smallcnn-architecture}
\end{figure}
